\section{An explicit source adjoint with interface conditions}\label{sec:source-adjoint}

\subsection{A traction objective on a fixed interface}
We now derive spatial optimality equations, rather than leaving the inverse
as an abstract residual. This section uses exactly the canonical transmission
model of \cref{sec:sources}. Its geometry is fixed during a source variation;
the coupled shape derivative is a separate operation. The pressure trace $P$
and normal-velocity trace $v$ are continuous across the interface. Tangential
velocity obeys $\bm V_{t\ell}=\Gs P/(\mathrm i\omega\rho_\ell)$.
Define three real material coefficients
\begin{equation}
 c_P=\frac14\left(\frac{1}{\rho_-c_-^2}-\frac{1}{\rho_+c_+^2}\right),
 \qquad c_v=\frac{\Delta\rho}{4},\qquad
 c_\tau=-\frac{1}{4\omega^2}\left(\frac{1}{\rho_-}-\frac{1}{\rho_+}\right).
 \label{eq:traction-trace-coefficients}
\end{equation}
Substitution into \cref{eq:radiation-flux,eq:acoustic-traction} gives
\begin{equation}
 \boxed{\Pi=c_P|P|^2+c_v|v|^2+c_\tau|\Gs P|^2.}
 \label{eq:traction-trace}
\end{equation}
All three terms have pressure units. The tangential radiation-traction jump
vanishes in this inviscid transmission limit, since $\rho_\ell\bm V_{t\ell}$
and $v$ are common on both sides. This is not a claim about thermoviscous
boundary layers or variable surface tension.

Let $\Pi_{\rm tar}$ be a required load in pascals, $A_0>0$ a reference area,
$\Pi_0>0$ a reference stress, and $\varpi\geq0$ a smooth dimensionless
weight supported away from the contact line. For a closed interface the last
restriction is unnecessary. Choose positive source-penalty weights
$\gamma_g,\gamma_F,\gamma_q$ for available components, using the scales of
\cref{eq:source-norm}, and define
\begin{align}
 J_{\rm load}&=\frac{1}{2A_0\Pi_0^2}
             \int_\Gamma\varpi(\Pi-\Pi_{\rm tar})^2\dd A,
             \label{eq:adjoint-load-cost}\\
 J_{\rm reg}&=\frac{\gamma_g}{2S_ag_0^2}\int_{\Gamma_a}|g|^2\dd A
       +\frac{\gamma_F}{2V_sF_0^2}\int_{D_s}|\bm F|^2\dd V
       +\frac{\gamma_q}{2V_sq_0^2}\int_{D_s}|q|^2\dd V.
       \label{eq:adjoint-source-penalty}
\end{align}
These costs are dimensionless. For a volume-constrained quiescent target use
$\Pi_{\rm tar}=\sigma\kappa+\Delta\rho\Phi-\lambda_V$ and optimize over
$\lambda_V$ as well. Its stationarity equation is the weighted removal of
the otherwise arbitrary constant load. A weighted patch objective does not
certify the unobserved part of the interface.

Define the real residual weight $r=\varpi(\Pi-\Pi_{\rm tar})/(A_0\Pi_0^2)$
and the complex interface covectors
\begin{equation}
 A_{\rm obs}=2c_PrP-2\Gs\cdot(c_\tau r\Gs P),
 \qquad B_{\rm obs}=2c_vrv.
 \label{eq:interface-covectors}
\end{equation}
Their units are $\mathrm{Pa^{-1}\,m^{-2}}$ and $\mathrm{s\,m^{-3}}$,
respectively. At fixed target and geometry, differentiating the three squared
moduli and integrating the tangential term by parts gives
\begin{equation}
 DJ_{\rm load}[\delta P,\delta v]
       =\Rea\int_\Gamma
          (\bar A_{\rm obs}\delta P+\bar B_{\rm obs}\delta v)\dd A.
 \label{eq:load-state-variation}
\end{equation}
For nonzero edge weights an additional term is
$2\Rea\int_{\mathcal C}c_\tau r(\bm m\cdot\Gs\bar P)\delta P\dd s$,
where $\bm m$ is the outward unit co-normal tangent to $\Gamma$. It vanishes
under the stated support restriction. Omitting it for a full-edge objective
would require a justified edge condition, not a change of notation.

\subsection{The adjoint has jumps even when pressure does not}
Let $Z_\ell$ be a piecewise adjoint field; for the normalization above its units
are $\mathrm{s^2\,m^{-3}}$, so it is not a physical pressure. Assume sufficient
piecewise regularity for the following Green identity and traces. With
$\partial_n$ retaining the same orientation on both sides of $\Gamma$, solve
\begin{align}
 L_\ell^*Z_\ell&=-\nabla\cdot(\alpha_\ell\nabla Z_\ell)
                           -\bar\beta_\ell Z_\ell=0
                  &&\text{in }\Omega^\ell,\label{eq:source-adjoint-bulk}\\
 \alpha\partial_{n_a}Z+\mathrm i\omega\bar Y Z&=0
                  &&\text{on }\partial\mathcal D_c,\label{eq:source-adjoint-wall}\\
 \boxed{\jump{Z}=\frac{\mathrm i}{\omega}B_{\rm obs},\qquad
 \jump{\alpha\partial_nZ}=-A_{\rm obs}}
                  &&\text{on }\Gamma.\label{eq:source-adjoint-jumps}
\end{align}
The bulk coefficient $\beta_\ell$ is real in the present canonical model;
the conjugation records the formal-adjoint convention. A model with complex
bulk coefficients also needs its consistently changed radiation and power
expressions. The plus sign in the adjoint wall impedance follows from the
chosen $e^{-\mathrm i\omega t}$ convention.

To verify the interface signs, write $\delta S=-\nabla\cdot
(\alpha\delta\bm F)-\mathrm i\omega\delta q$, the volume right-hand side
of the differentiated wave equation. Green's second identity, summed over
the phases, gives
\begin{align}
 \int_{D_s}\delta S\bar Z\dd V
   &=\int_\Gamma\left[-\delta P\jump{\alpha\partial_n\bar Z}
                      +\mathrm i\omega\delta v\jump{\bar Z}\right]\dd A
         -\mathrm i\omega\int_{\Gamma_a}\delta g\bar Z\dd A.
 \label{eq:source-green-identity}
\end{align}
The pressure variation is continuous and its common flux is
$\alpha\partial_n\delta P=\mathrm i\omega\delta v$.
The two boxed jumps turn the interface integral into the complex expression
whose real part is \cref{eq:load-state-variation}. Integrating the body-force
source by parts, using its interior support, therefore yields
\begin{equation}
 DJ_{\rm load}[\delta u]=\Rea\left[
  \mathrm i\omega\int_{\Gamma_a}\delta g\bar Z\dd A
  +\int_{D_s}\alpha\delta\bm F\cdot\nabla\bar Z\dd V
  -\mathrm i\omega\int_{D_s}\delta q\bar Z\dd V\right].
 \label{eq:explicit-source-variation}
\end{equation}
This derivation explains why an adjoint constrained to globally continuous
$Z$ would generally be wrong for an objective depending on normal velocity.
An equivalent mixed or transposition formulation must encode these same
interface functionals. Boundary terms are part of the derivative, as in
continuous-adjoint methodology \cite{giles2000}.

\subsection{Spatial source gradients and necessary optimality}
Define gradients using the unweighted real physical $L^2$ pairing:
\begin{equation}
 D(J_{\rm load}+J_{\rm reg})[\delta u]
  =\Rea\left[\int_{\Gamma_a}\bar G_g\delta g\dd A
        +\int_{D_s}\bar{\bm G}_F\cdot\delta\bm F\dd V
        +\int_{D_s}\bar G_q\delta q\dd V\right].
 \label{eq:physical-gradient-pairing}
\end{equation}
Combining \cref{eq:adjoint-source-penalty,eq:explicit-source-variation} gives
\begin{equation}
 \boxed{G_g=\frac{\gamma_g g}{S_ag_0^2}-\mathrm i\omega Z,\qquad
 \bm G_F=\frac{\gamma_F\bm F}{V_sF_0^2}+\alpha\nabla Z,\qquad
 G_q=\frac{\gamma_q q}{V_sq_0^2}+\mathrm i\omega Z.}
 \label{eq:distributed-gradients}
\end{equation}
Each field is restricted to its own support. The gradient in the scaled
Hilbert metric of \cref{eq:source-norm} has the corresponding metric factors;
the physical first variation is the same. If $\lambda_V$ is free, also impose
$\int_\Gamma r\dd A=0$.

For unconstrained boundary actuation, the first-order condition is the
implicit field equation
\begin{equation}
 g=\frac{\mathrm i\omega S_ag_0^2}{\gamma_g}Z|_{\Gamma_a}.
 \label{eq:boundary-source-optimality}
\end{equation}
For the sole additional restriction $|g(\bm x)|\leq g_{\max}(\bm x)$,
where $g_{\max}\geq0$ is prescribed, this becomes pointwise projection of
the right-hand side onto the complex disk of that radius. Phase-only control,
source coupling or a spatial-derivative penalty does not have this simple
projection rule. In a general convex admissible set the condition is that
\cref{eq:physical-gradient-pairing}, with $\delta u=\tilde u-u$, is
nonnegative for every $\tilde u\in\mathcal U_{\rm ad}$.

These equations form a common state--adjoint--optimality system for every
admissible target. They are not a closed-form answer: $Z$ depends on the
current pressure, load residual and geometry. The quadratic loading makes
the objective nonconvex, so the system describes necessary local conditions.
A vanishing source is always a stationary point for a pure quadratic-load
tracking objective with quadratic regularization; a gradient-only iteration
started exactly there can stop even when a useful nonzero source exists.
That analytic degeneracy must be addressed by the eventual solution method,
not concealed by a target-specific fit.

\subsection{A limiting boundary check and extension to the full model}
For a one-sided, fixed pressure-release surface, $P=0$ and
$\Pi=\rho_-|v|^2/4$. Let $B_{\rm obs}$ now be the corresponding velocity
covector in \cref{eq:load-state-variation}. The exterior Green boundary term
is $-\mathrm i\omega\delta v\bar Z$, so the adjoint Dirichlet condition is
\begin{equation}
 Z=-\frac{\mathrm i}{\omega}B_{\rm obs}\quad\text{on }\Gamma.
 \label{eq:pressure-release-adjoint}
\end{equation}
Setting $Z=0$ there would lose the entire velocity-based objective derivative.
This limit is a sign and boundary-condition check, not the assumed apparatus.

For the full coupled problem, the same integration-by-parts construction must
include viscous and thermal fields, fast interface compliance, the mean flow,
actuator loading and geometric transport. \Cref{eq:distributed-gradients}
is complete for the explicitly stated frozen-geometry acoustic subproblem;
it is not the full shape or thermoviscous adjoint. The general coupled and
time-dependent adjoint equations are specified in \cref{sec:control}.
